%% normalization.tex -- shared sub-recipe. Referenced by every analysis that
%% consumes expression vectors, via \dependson{sub:normalization}.

\begin{recipe}[label=sub:normalization]{Normalization of Expression Values}

\recipepurpose{Turn a replicate-level expression table into the gene
  $\times$ tissue matrix that every Tensor Omics analysis uses as its
  expression vectors: gene-wise scale removed, replicates averaged per tissue,
  dynamic range compressed. This is the shared entry point of the TOX
  pipeline; the geometry of every later recipe is fixed by the choices made
  here.}

\nomaintainer

\begin{requiredinputs}
  \item A table of non-negative expression measurements (e.g.\ TPM from
        \cmd{kallisto}), genes $\times$ replicates, one column per biological
        replicate.
  \item The replicate layout: how many replicates belong to each tissue,
        condition or time point. The replicates of one tissue must occupy
        \emph{adjacent} columns.
  \item For stress-response designs only: which tissue is the control and
        which the condition, for each pair you want a fold change of.
\end{requiredinputs}

\begin{prerequisites}
  \item None --- this is the first step of every Tensor Omics analysis.
\end{prerequisites}

\begin{workflow}
  \step{Arrange the measurements as an \code{(n\_replicates, n\_genes)}
        column-major array, so that \emph{one gene is one column}. Group the
        replicates of each tissue into adjacent rows and record the group
        sizes in \code{reps\_per\_tissue}, e.g.\ \code{[4, 5, 5]} for three
        tissues measured with four, five and five replicates.}
  \step{Decide whether to quantile-normalize across replicates
        (\optn{use\_quantile}). It is off by default; switch it on when the
        replicate distributions differ enough that the difference is technical
        rather than biological.}
  \step{Run \code{normalization\_pipeline}. It scales each gene by a
        LOESS-stabilized standard deviation, optionally quantile-normalizes,
        averages the replicates of each tissue, and finally applies
        $x \mapsto \log_2(x+1)$.}
  \step{Check the result against the expectations below before using it:
        the shape, the value range, and the number of genes that reached the
        output unscaled.}
  \step{Stress-response designs only: call \code{calc\_fchange} on the
        \emph{log-transformed} matrix to turn control/condition pairs into
        log2 fold changes, so each such axis carries e.g.\ ``drought in root
        over control root'' rather than an absolute level.}
\end{workflow}

\examplescript{python/examples/normalization.py}{%
  Reads a replicate-level TSV, derives the tissue layout from its header,
  runs the pipeline and writes the normalized matrix. Run unchanged from the
  repository root, it normalizes the bundled
  \file{material/kallisto_sex_data_no_na.tsv} (88\,327 genes, 67 replicates,
  13 tissues) in a few seconds and writes
  \file{results/normalized_expression.tsv}.}

\begin{codefragment}[language=Python,caption={Fragments to edit: the input, the tissue layout, the parameters}]
import numpy as np
from tensor_omics import normalization_pipeline

# --- input: one gene per COLUMN, --------------------------
#     replicates of a tissue in adjacent ROWS
expr = np.asfortranarray(raw_counts)  # (n_replicates, n_genes)

# --- tissue layout: one entry per tissue, -----------------
#     summing to n_replicates
reps_per_tissue = np.array([4, 5, 5], dtype=np.int32)

# --- parameters (see the table below) ---------------------
normalized = normalization_pipeline(
    expr, reps_per_tissue,
    span=0.7,            # LOESS span, mean-vs-SD fit
    degree=2,            # LOESS polynomial degree
    use_quantile=False,  # quantile normalization
)                        # -> (n_tissues, n_genes)
\end{codefragment}

\begin{note}
  Tensor Omics stores one expression vector per \emph{column} because Fortran
  is column-major. A matrix handed over as genes $\times$ replicates is not
  rejected --- it is silently interpreted as replicates $\times$ genes, and
  every downstream distance and angle is then meaningless. Check
  \code{expr.shape} before you call.
\end{note}

\begin{parameters}
  \param{reps\_per\_tissue}{How the replicate rows are grouped into tissues;
    entry $i$ is the number of adjacent rows belonging to tissue $i$}{%
    Must sum to \code{n\_replicates}, otherwise the call fails with
    \code{Array size mismatch}. Replicate counts may differ between tissues}
  \param{span}{Fraction of the genes entering each local fit of the
    mean--SD LOESS curve}{Keep the default \code{0.7}.
    Lower it only if the mean--SD trend is genuinely non-monotonic and the
    fit visibly oversmooths; too low a span makes the curve follow noise}
  \param{degree}{Polynomial degree of the local LOESS fits}{%
    Keep the default \code{2}. Use \code{1} for a strongly linear trend or
    very few genes}
  \param{use\_quantile}{Whether the replicate distributions are forced onto a
    common distribution after gene-wise scaling}{Default \code{False}. Switch
    on when replicates differ in distribution for technical reasons; leave off
    when the between-replicate differences are the signal}
\end{parameters}

\begin{expectedresults}
  The result has shape \code{(n\_tissues, n\_genes)} and is read-only ---
  call \code{.copy()} if you need to modify it. For a non-negative input every
  value is $\geq 0$, since $\log_2(x+1) \geq 0$ there. On the bundled example
  data the values span $[0, 4.06]$ without quantile normalization and
  $[0, 4.27]$ with it, and 2\,627 of the 88\,327 genes (3\%) have zero
  variance across their replicates and therefore reach the output unscaled.
  A value range wildly different from this, or a much larger unscaled
  fraction, points at the input layout rather than at biology.
\end{expectedresults}

\begin{pitfall}[Pitfall: replicates that are not adjacent]
  \code{reps\_per\_tissue} describes \emph{contiguous blocks} of rows, not a
  grouping by name. If the replicate columns of a tissue are scattered through
  the table, the pipeline still runs and silently averages the wrong columns
  together. Sort the columns tissue by tissue first; the example script
  refuses to continue when it detects a non-adjacent tissue.
\end{pitfall}

\begin{pitfall}[Pitfall: genes the standard-deviation fit cannot use]
  The scaling factor comes from a LOESS curve fitted over all genes, so a gene
  with zero variance across its replicates carries no mean-versus-SD
  information. Such genes are left out of the fit \emph{and reach the output
  unscaled} --- they are not dropped and not flagged. Count them, as the
  example script does. If fewer than five genes have non-zero variance, the
  whole call fails with \code{ToxInputError: invalid input arguments}; this is
  what you get for an all-zero matrix, or for a toy example with too few
  genes.
\end{pitfall}

\begin{pitfall}[Pitfall: expecting quantile normalization by default]
  Quantile normalization is an \emph{option}, not a pipeline stage: the
  default is scale $\rightarrow$ average $\rightarrow$ $\log_2(x+1)$. Two
  analyses that differ only in \optn{use\_quantile} live in different
  geometries and must not be compared. Record the flag alongside the results.
\end{pitfall}

\begin{pitfall}[Pitfall: fold changes computed on the wrong matrix]
  \code{calc\_fchange} \emph{subtracts} one axis from another. That is a log2
  fold change only if you pass the log-transformed matrix --- which is what
  \code{normalization\_pipeline} returns. Passing raw tissue averages instead
  yields a plain difference of expression levels that merely looks like a fold
  change.
\end{pitfall}

\begin{warning}
  \code{root\_mean\_sq\_normalization} divides by
  $\sqrt{\mathrm{mean}(x^2)}$, which is not a standard deviation. It is kept
  for possible future use and must not be used for normalization; use
  \code{normalization\_pipeline} or \code{normalize\_by\_std\_dev}.
\end{warning}

\begin{interpretation}
  After normalization, a gene's coordinates state \emph{where its expression
  sits across tissues}, not how many transcripts were counted. Gene-wise
  scaling removes each gene's own expression magnitude, so two genes differing
  a thousandfold in abundance can be neighbours if their tissue profiles agree;
  the $\log_2(x+1)$ step then compresses the remaining dynamic range so that no
  single highly expressed tissue dominates a distance. Distances and angles in
  this space therefore compare the \emph{shape} of expression across tissues.

  Because each choice made above defines a different space, results obtained
  under different settings answer different questions and generally disagree.
  Agreement between them is evidence of a robust effect; disagreement is
  usually informative rather than an error. Report the normalization settings
  with any result derived from this matrix.
\end{interpretation}

\begin{references}
  \reference{Cleveland (1979). Robust locally weighted regression and
    smoothing scatterplots. \emph{Journal of the American Statistical
    Association} 74(368):829--836.}
  \reference{Cleveland \& Devlin (1988). Locally weighted regression: an
    approach to regression analysis by local fitting. \emph{Journal of the
    American Statistical Association} 83(403):596--610.}
  \reference{Bolstad, Irizarry, {\AA}strand \& Speed (2003). A comparison of
    normalization methods for high density oligonucleotide array data based on
    variance and bias. \emph{Bioinformatics} 19(2):185--193.}
\end{references}

\end{recipe}
